\documentclass[../../../include/open-logic-section]{subfiles}

\begin{document}

\olfileid{sth}{spine}{recursion}
\olsection{مبرهنة (مبرهنات) العودية المتناهية التجاوز}

أول ما يجب فعله هو التحقق من نجاح~\olref[valpha]{defValphas} بوصفه
تعريفًا. وبوجه أعم، يلزمنا أن نثبت نجاح أي محاولة لتقديم تعريف
بالعودية (المتناهية التجاوز). وهذه غاية هذا القسم.

\emph{تحذير: هذه مادة دقيقة}. لكن المغزى العام بسيط: يضمن الاستقراء
المتناهي التجاوز مع الاستبدال مشروعية (صور عدة من) العودية المتناهية
التجاوز.\footnote{تذكير: يجوز أن تكون لجميع الصيغ والحدود معاملات
  (ما لم يُصرح بخلاف ذلك).}

\begin{defn}
	ليكن~$\tau(x)$ حدًا، ولتكن~$f$ دالة، وليكن~$\alpha$ عددًا ترتيبيًا.
	نقول إن~$f$ \emph{تقريب~$\alpha$} لـ$\tau$ إذا وفقط إذا كان كل من
	$\dom{f} = \alpha$ و
	$(\forall \beta \in \alpha)f(\beta) =
	\tau(\funrestrictionto{f}{\beta})$.
\end{defn}

\begin{lem}[العودية المحدودة]\ollabel{transrecursionfun}
لكل حد~$\tau(x)$ ولكل عدد ترتيبي~$\alpha$، يوجد تقريب~$\alpha$ وحيد
لـ$\tau$.
\end{lem}
\begin{proof}
سنبيّن أنه لكل~$\gamma \leq \alpha$ يوجد تقريب~$\gamma$ وحيد.

نثبت الوحدانية أولًا. لتكن~$g$ و$h$، على الترتيب، تقريبي~$\gamma$
و$\delta$. يبيّن استقراء متناهٍ في التجاوز على وسيطاتهما أن
$g(\beta) = h(\beta)$ لكل
$\beta \in \dom{g} \cap \dom{h} =
\gamma \cap \delta = \min(\gamma, \delta)$. ولذلك تكون تقريباتنا وحيدة
(إذا وجدت)، وتتفق في جميع القيم.

ولإثبات الوجود، نستعمل الآن استقراءً بسيطًا متناهيًا في التجاوز
(\olref[ordinals][opps]{simpletransrecursion}) على الأعداد الترتيبية
$\delta \leq \alpha$.

الدالة الخالية تقريب~$\emptyset$ على نحو بديهي.

إذا كانت~$g$ تقريب~$\gamma$، فإن
$g \cup \{\tuple{\gamma, \tau(g)}\}$ تقريب~$\ordsucc{\gamma}$.

إذا كان~$\gamma$ عددًا ترتيبيًا حديًا، وكانت~$g_\delta$ تقريب~$\delta$
لكل~$\delta < \gamma$، فلتكن
$g = \bigcup_{\delta \in \gamma} g_\delta$. وهذه دالة لأن مختلف
دوال~$g_\delta$ لدينا تتفق في جميع القيم. وإذا كان
$\delta \in \gamma$، فإن
$g(\delta) = g_{\ordsucc{\delta}}(\delta) =
\tau(\funrestrictionto{g_{\ordsucc{\delta}}}{\delta}) =
\tau(\funrestrictionto{g}{\delta})$.

يتم بذلك البرهان بالاستقراء المتناهي التجاوز.
\end{proof}

إذا سمحنا لأنفسنا بتعريف \emph{حد} بدلًا من دالة، استطعنا حذف القيد
$\alpha$ من النتيجة السابقة. وفي تقرير النتيجة الآتية وبرهانها، عندما
تكون~$\sigma$ حدًا، نجعل
$\funrestrictionto{\sigma}{\alpha} = \Setabs{\tuple{\beta,
\sigma(\beta)}}{\beta \in \alpha}$.

\begin{thm}[العودية العامة]\ollabel{transrecursionschema}
لكل حد~$\tau(x)$، نستطيع أن نعرّف صراحة حدًا~$\sigma(x)$ بحيث
$\sigma(\alpha) = \tau(\funrestrictionto{\sigma}{\alpha})$ لكل عدد
ترتيبي~$\alpha$.
\end{thm}

\begin{proof}
لكل~$\alpha$، يوجد بموجب~\olref{transrecursionfun} تقريب~$\alpha$
وحيد~$f_\alpha$ لـ$\tau$. عرّف~$\sigma(\alpha)$ بأنه
$f_{\ordsucc{\alpha}}(\alpha)$. والآن:
	\begin{align*}
		\sigma(\alpha) &= 
		f_{\ordsucc{\alpha}}(\alpha) \\&= 
		\tau(\funrestrictionto{f_{\ordsucc{\alpha}}}{\alpha}) \\&= 
		\tau(\Setabs{\tuple{\beta, f_{\ordsucc{\alpha}}(\beta)}}{\beta \in \alpha}) \\&= 
		\tau(\Setabs{\tuple{\beta, f_{\ordsucc{\beta}}(\beta)}}{\beta \in \alpha})\\&=
		\tau(\funrestrictionto{\sigma}{\alpha})
	\end{align*}
مع ملاحظة أن
$f_{\ordsucc{\beta}}(\beta) = f_{\ordsucc{\alpha}}(\beta)$ لكل
$\beta < \alpha$، كما في~\olref{transrecursionfun}.
\end{proof}
\noindent
لاحظ أن~\olref{transrecursionschema} \emph{مخطط}. والأهم أننا لا
نستطيع توقع أن تعرّف~$\sigma$ دالة، أي نوعًا معينًا من
\emph{المجموعات}، لأن~$\dom{\sigma}$ ستكون عندئذ مجموعة جميع الأعداد
الترتيبية، خلافًا لمفارقة بورالي--فورتي
(\olref[ordinals][basic]{buraliforti}).

يبقى أن نبيّن أن~\olref{transrecursionschema} يسوّغ تعريفنا لمجموعات
$V_\alpha$. وقد لا يكون ذلك واضحًا فورًا؛ لكنه سيظهر مع صورة أخيرة
بسيطة من العودية المتناهية التجاوز.

\begin{thm}[العودية البسيطة]\ollabel{simplerecursionschema}
لكل حدين~$\tau(x)$ و$\theta(x)$ ولكل مجموعة~$A$، نستطيع أن نعرّف
صراحة حدًا~$\sigma(x)$ بحيث:
\begin{align*}
	\sigma(\emptyset) &= A\\
	\sigma(\ordsucc{\alpha}) &= \tau(\sigma(\alpha)) &&
		\text{لكل عدد ترتيبي }\alpha\\
	\sigma(\alpha) &= \theta(\ran{\funrestrictionto{\sigma}{\alpha}})&&
		\text{عندما يكون }\alpha\text{ عددًا ترتيبيًا حديًا}
\end{align*}
\end{thm}

\begin{proof}
نبدأ بتعريف حد~$\xi(x)$ كما يأتي:
%t $\xi(x) = A$ if $x$ is not a function whose domain is an ordinal;
%otherwise let:
\[
	\xi(x) = 
	\begin{cases}
			A &\text{إذا كان $x = \emptyset$، أو لم يكن $x$ دالة}\\
			&\hspace{1em}\text{مجالها عدد ترتيبي؛ وخلاف ذلك:}\\
			\tau(x(\alpha)) & \text{إذا كان $\dom{x} = \ordsucc{\alpha}$}\\
			\theta(\ran{x}) & \text{إذا كان $\dom{x}$ عددًا ترتيبيًا حديًا}
	\end{cases}
\]
بموجب~\olref{transrecursionschema}، يوجد حد~$\sigma(x)$ يحقق
$\sigma(\alpha) = \xi(\funrestrictionto{\sigma}{\alpha})$ لكل عدد
ترتيبي~$\alpha$؛ وفوق ذلك يكون~$\funrestrictionto{\sigma}{\alpha}$
دالة مجالها~$\alpha$. ونبيّن أن~$\sigma$ لها الخصائص المطلوبة
بالاستقراء البسيط المتناهي التجاوز
(\olref[ordinals][opps]{simpletransrecursion}).

أولًا، $\sigma(\emptyset) = \xi(\emptyset) = A$.

ثم
$\sigma(\ordsucc{\alpha}) =
\xi(\funrestrictionto{\sigma}{\ordsucc{\alpha}}) =
\tau(\funrestrictionto{\sigma}{\ordsucc{\alpha}}(\alpha)) =
\tau(\sigma(\alpha))$.

وأخيرًا،
$\sigma(\alpha) = \xi(\funrestrictionto{\sigma}{\alpha}) =
\theta(\ran{\funrestrictionto{\sigma}{\alpha}})$، عندما يكون~$\alpha$
حديًا.
\end{proof}
\noindent
والآن، لتسويغ~\olref[valpha]{defValphas}، خذ ببساطة
$A = \emptyset$ و$\tau(x) = \Pow{x}$ و$\theta(x) = \bigcup x$.
وهكذا، بعد طول انتظار، يسوَّغ تعريف مجموعات~$V_\alpha$!{}


\end{document}
